\chapter{撰写研究经费申请书}
\chapterauthors{Steven R. Cummings、Deborah G. Grady 和 Alka M. Kanaya}
\sourcepage{615}
如第2章所述，研究开发始于陈述问题，随后写一页计划提纲（附录1），请导师、同事和专家建议。在临床研究设计课中，我们要求学生写5--7页计划，包括背景、设计、参与者、预测与结局及潜在混杂变量测量、样本量和效能估计，以及人体研究问题。

\term{研究方案}是详细书面计划，迫使研究者组织、澄清、完善各环节，提高科学方法的周密程度与效率。即使不需资助，仍需方案指导工作并获得伦理批准。\term{申请书}（proposal）则用于向资助机构请求经费，包括目标、意义、方法、人体研究问题，以及预算、行政和支持信息。

本章介绍如何写出更可能获资助的申请，重点是 NIH 格式的原创研究；退伍军人事务部、CDC、国防部、医疗保健研究与质量局及私人基金会等，多采用相似格式。NIH 网站提供写作、预算和提交建议：\url{https://grants.nih.gov/grants/how-to-apply-application-guide.html}。
\sourcepage{616}
\section{撰写申请书}
准备通常需要数月组织、写作和修订，以下步骤有助起步。
\begin{itemize}
\item \textbf{决定投向。}各机构兴趣、程序、要求不同，应先决定申请对象、金额上限，取得特定写法和截止期指南。可从 NIH（\url{http://grants.nih.gov/grants/oer.htm}）开始。各研究所及其他资助方网站列出优先方向；也可联系项目官员，其联系方式和职责见资助机会公告及网站。
\item \textbf{组建团队、指定负责人。}多数申请由日后实施团队共同撰写，可小至研究者和导师，也可包括合作者、生物统计师、财务人员、助理和辅助人员；必须具备或能取得设计实施所需的主要专长。

PI 领导全程，对研究具有最终权力并承担责任；应持续带领团队，分派写作和其他任务、设时限、定期开会、确保按期完成并负责质量。PI 常是资深科学家，其知识判断帮助设计，既往成绩提高成功和获资助的可能。
\sourcepage{617}
不过，NIH 鼓励新研究者作 PI，设有专门机会，对职业早期申请常给予优先（\url{http://grants.nih.gov/grants/new_investigators/}）。NIH 将尚未担任其研究经费 PI 者称为\term{新研究者}（new investigators），在最终研究学位或临床训练完成后10年内者为\term{职业早期研究者}（early stage investigators）。初次 PI 若已在资深科学家指导和经费、职业发展奖励或机构基金会小额资助下开展过研究，较易获批。原创论文第一作者等发表记录，显示独立潜力和领导准备。

首次 PI 应纳入在相关领域有成功经历的共同研究者，指导实施、改善评审与获资助机会。通常一份申请一名 PI，但 NIH 允许专长互补、角色职责明确的多 PI（\url{http://grants.nih.gov/grants/multi_pi/overview.htm}）。
\item \textbf{遵循指南。}写前阅读资助类型、组织格式、页数、金额等书面要求。但指南不能涵盖机构运行和偏好的一切，早期应与其人员讨论兴趣、范围和详略。
\sourcepage{618}
NIH、其他联邦机构及基金会常有科学管理人员，即\term{项目官员}（program officers），帮助申请回应优先方向。可电邮或电话澄清指南、兴趣和评审程序，再在会议或到访时建立工作关系。

列出必需细节，提交前反复核查。优秀申请也可能因不合规定被拒，令人沮丧却可避免。大学经费管理员虽常有检查清单，遵循规定最终仍是 PI 责任。
\item \textbf{制定时间表、定期开会。}书面任务时限能温和督促履责。除科学内容，还需考虑本机构行政要求：场地、预算、分包审查可能耗时，真正完成期限可能比资助方早几天至几周。拖到最后会造成危机，损害本来良好的申请。时间表宜规定各书面产物时限，让人人参与设定任务；会议或电话中检查进度及时间是否现实。
\item \textbf{寻找范例。}查看目标机构近期资助申请，了解良好内容和格式，启发想法，帮助写得更清楚、合理、有说服力。
\sourcepage{619}
也应审阅成功或失败申请的书面批评，了解评审者重视什么；可从机构科研资助办公室或类似研究 PI 处获取。
\item \textbf{从提纲开始。}先按表20.1组织任务，多人写作时注明各部分负责人。不要以为完整研究计划必须全部想好才能写第一句；先把想法写下，作为编辑、完善和具体请教的材料。
\item \textbf{反复审阅修订。}申请常有许多版本，纳入新想法、建议和数据。尽早请熟悉主题和资助机构的同事批评，着重意义、创新、设计方法效度、组织和清晰度。提交前尖锐详尽的批评，好过未预见问题而被拒。接近完成时检查内部一致性、指南、语法、格式和错字；草率写作会让人怀疑工作和领导能力，损害好想法。
\end{itemize}

\begin{bookbox}{表20.1\quad 以 NIH 为例的申请书主要要素}
\begin{description}[leftmargin=1em,style=nextline]
\item[标题]
\item[项目概要或摘要]
\sourcepage{620}
\item[行政部分]预算与理由；研究者个人简历；机构资源。
\item[具体目标（通常1页）]
\item[潜在影响]
\item[研究策略（通常12页）]意义；创新；研究方法，包括：
\begin{itemize}
\item 总体概述；
\item 依据（理由和初步数据）；
\item 参与者：选择标准、抽样设计、招募计划、优化依从性和完整随访的计划；
\item 适用的程序：随机化、盲法；
\item 测量：主要预测变量（试验为干预）、结局变量、潜在混杂变量；
\item 统计：分析方法；假设、样本量和效能；
\item 访视内容与时间；数据管理和质控；时间表与组织图；局限和其他方法。
\end{itemize}
\item[人体研究]
\item[参考文献]
\item[附录与合作协议]
\end{description}
\small NIH：美国国立卫生研究院。
\end{bookbox}

\section{大型经费申请的组成}
表20.1列出 NIH R01 等大项目要素。其他 NIH 经费合同及其他机构相似，但信息可较少或格式不同，应查具体指南。
\sourcepage{621}
\subsection{开头部分}
描述性、简洁标题给人关于目标和设计的第一印象，也便于日后记住。例如，“MRI引导高频超声与假超声治疗有症状子宫肌瘤的随机试验”概括问题、设计和人群。避免“旨在确定……的一项研究”等多余措辞。

摘要先介绍目标和依据，再述设计方法，最后说明潜在发现的影响。既要对本领域和相关人员有信息，也要让有科学素养的普通读者理解。多数机构限制字数或行数，措辞应精练而有描述性。有些评审者只读摘要；对其他人，它也是便于回忆细节的概述。因此必须独立成立，涵盖所有主要特点，有说服力地描述优势和影响。

\subsection{行政部分}
几乎所有机构要求预算、理由和机构资源，说明人员资格、机构条件、设备场地与专长的获取。预算按指南组织。NIH 要求首12个月详细预算和整个项目（通常2--5年）汇总预算。
\sourcepage{622}
详细预算包括：人员（姓名、职位、投入时间比例、工资福利）；咨询；设备；耗材；差旅；患者照护；改建装修；联合／合同；其他如电话、邮资、会议通话、复印、插图、发表、书籍和付费服务。

不要最后才做预算，场地设备人员估价需时间。大学行政人员熟悉预算、理由和其他行政内容，应尽早通知计划，定期开会检查进展时限。提纲一成即可开始，建议金额、防止漏项，预见机构规则和延误，并协助撰写理由、资源说明、收集简历、分包、附录等。

每项金额都需充分说明。人员通常占大部分成本，应以需要和具体职责论证投入比例；完整而简洁的职位说明，应让评审确信每人投入都必不可少。评审常关注关键人员时间：过少且其他任务多，暗示无力投入；更常见的是认为比例高于职位所需。
\sourcepage{623}
若理由不足，评审可建议削减或删除。

再周全的预算实际也会改变。获批后，若变动适度且用于研究目标，通常可改变花费方式；跨类别调款或关键研究者投入改变超过25\%，可能需先批准。只要不增加总额，合理重编常获同意。

NIH 要求所有获酬研究者和顾问提交\term{个人简历}（biosketch），按指定五页格式，含经验与研究相关性的个人陈述、教育培训任职、有限数量的论文荣誉，以及相关经费合同。

机构资源可包括计算机技术设备、特殊影像测量、办公室实验室、招募、收集管理及标本保存条件，常借用既往申请或机构、中心、实验室的标准描述。

\subsection{具体目标}
\term{具体目标}（specific aims）以具体语言陈述问题和期望成果，NIH 限1页。\textbf{很多评审最重视此页，应仔细写作，并随申请发展反复修订。}
\sourcepage{624}
常先用两三短段概述问题及重要性、已有研究为何未解决、拟采取方法，再简述可实际完成的描述性目标，尽可能表达为可检验假设。

目标按适合研究的逻辑排列，可先横断面再随访，先病理生理机制再临床或公共卫生结局。职业发展申请常先定性，用焦点小组开发工具或干预，再定量讨论预测、结局和检验，即混合方法。也可把最重要目标置首，使其在样本量效能等各部分都优先出现。

最后通常短述对健康疾病知识、实践、公共卫生或未来研究的潜在影响，旨在使包括可能只读此页的非主次评审在内，全体委员给予积极评价。

\subsection{研究策略}
原书所述 NIH 格式将多数申请的研究策略限12页，分三部分：
\begin{itemize}
\item \term{意义}（significance），通常2--3页，说明发现如何增进科学认识、解决重要问题或领域障碍、改善实践与公共卫生或影响政策。介绍问题规模、已有成果、知识缺口，以及本研究如何推动进展。
\sourcepage{625}
\item \term{创新}（innovation），通常1--2页，指出新测量、疾病机制、不同或更大人群、新治疗、分析方法等进步。NIH 关注如何以创新概念、方法和干预改变当前范式；但许多获资助项目只是渐进改善。建议准确描述新颖性，不夸称改变范式或全新方法。
\item \term{研究方法}（approach，过去称“methods”），通常7--9页，详述设计实施，受到最仔细审阅。NIH 建议按具体目标组织，包含表20.1内容和大致顺序。先以简明总体概述，有时配图表，使读者了解设计、参与者、主要测量、干预、随访长度和结局（表20.2）。
\end{itemize}

\begin{bookbox}{表20.2\quad 睾酮对心脏病、前列腺癌和骨折危险因素影响随机试验的访视安排}
\small\begin{tabularx}{\linewidth}{Yccccc}\toprule
 & 筛选 & 随机化 & 3个月 & 6个月 & 12个月\\\midrule
\sourcepage{626}病史 &X&--&--&--&X\\
血压 &X&X&X&X&X\\
前列腺检查 &X&--&--&--&X\\
前列腺特异抗原 &X&--&--&--&X\\
血脂 &--&X&X&X&X\\
炎症标志物 &--&X&--&--&X\\
骨密度 &--&X&--&--&X\\
骨转换标志物 &--&X&X&--&X\\
握力 &--&X&X&X&X\\
不良事件 &--&--&X&X&X\\\bottomrule
\end{tabularx}
\end{bookbox}

方法部分通常包括简短理由及\term{初步数据}（preliminary data），即团队既往研究提示拟议项目能够成功。强调既往工作重要性及为何继续扩展。先导结果对许多申请重要，尤其团队缺少方法经验、程序或招募可行性存疑时。\textbf{这是展示团队必要经验专长的机会。}
\sourcepage{627}
其他部分前文已讨论。参与者部分（第3章）界定并论证纳入排除、抽样、招募随访，使评审相信能达到人数，并说明优化干预与访视依从性。还需描述随机化盲法等重要程序，说明预测、结局、混杂变量如何及何时测量，干预如何实施、主要结局如何判定（第4章）。

统计部分通常按目标先述分析计划，例如先描述性汇总，再分析变量关联。随后讨论样本量和效能（第5、6章），从决定样本量的目标之零假设开始。效应量和测量精密度的假定须引用文献或初步工作；可用图表展示不同假定、效能和效应量如何影响样本量，说明选择合理。多数 NIH 小组很重视统计，应请统计师参与写作，复杂时列为研究者。
\sourcepage{628}
用表列访视、联系、时间和每次程序测量，可简明展示全部活动（表20.2）。质控与数据管理（第18、19章）应说明收集、储存、编辑，以及提高质量和安全的计划。

申请必须提供现实工作计划及各阶段开始、完成日期（图20.1），人员配置等也可用时间表；大型研究组织图还应显示权责、报告关系和运行方式。

\begin{figure}[htbp]\centering
\begin{tikzpicture}[x=.018cm,y=.018cm,font=\fontsize{8}{10}\selectfont]
\draw (0,0) rectangle (778,-263);
\foreach \y in {-29,-54,-96,-136,-177,-217} {\draw (0,\y)--(778,\y);}
\draw (120,0)--(120,-263);
\foreach \x in {275,418,562,707} {\draw (\x,0)--(\x,-54);\draw[dashed] (\x,-54)--(\x,-263);}
\foreach \x in {160,203,239,311,348,382,454,490,526,598,634,670,742} {\draw (\x,-29)--(\x,-54);}
\foreach \x in {203,348,490,634} {\draw[dashed] (\x,-54)--(\x,-263);}
\node at (60,-27){任务};
\node at (197,-14){第1年};\node at (346,-14){第2年};\node at (490,-14){第3年};\node at (634,-14){第4年};\node at (742,-14){第5年};
\foreach \x/\q in {140/1,181/2,221/3,257/4,293/1,329/2,365/3,400/4,436/1,472/2,508/3,544/4,580/1,616/2,652/3,688/4,724/1,760/2} {\node at (\x,-42){Q\q};}
\node[align=center,text width=2cm] at (60,-75){1. 准备测量工具};
\node[align=center,text width=2cm] at (60,-116){2. 招募参与者};
\node[align=center,text width=2cm] at (60,-156){3. 随访访视与资料收集};
\node[align=center,text width=2cm] at (60,-197){4. 清理数据};
\node[align=center,text width=2cm] at (60,-240){5. 分析与写作};
\draw[fill=white] (130,-62) rectangle (203,-92);
\draw[fill=gray!20] (183,-103) rectangle (260,-133);
\draw[fill=gray!35] (223,-146) rectangle (447,-175);
\draw[fill=gray!55] (403,-186) rectangle (550,-215);
\draw[fill=gray!70] (466,-227) rectangle (754,-257);
\end{tikzpicture}
\caption{准备申请书的假想时间表。（Q表示季度。）}
\end{figure}

虽非必需，讨论局限及替代方法常有帮助。不要忽略潜在缺点，而应说明权衡。指出重要困难及解决办法，可把批评转化为优点；但过度强调又会让评审过分关注弱处。目标是让其相信已预见主要问题，并有现实、深思熟虑的应对。
\sourcepage{629}
\subsection{大型申请的最后组成}
人体研究部分讨论知情同意、安全、隐私、保密（第7章），NIH 还要求说明风险收益告知与同意程序、纳排标准，以及全生命周期、女性和少数群体纳入计划与排除理由。试验另需详细招募留存、时间表、数据安全监察，并分别说明干预、设盲、分配、结局、统计设计效能和传播。虽不限页数，仍应简洁。

书目应全面、更新而有选择，体现熟悉领域，不应罗列一切。引文必须准确；错误引用或误解会使熟悉领域的评审反感。

有些申请可用附录放技术和支持细节，但 NIH 为防绕过页限，只允许收集工具（如问卷）和空白同意／赞同表。只有主要、次要评审收到附录，所以重要内容须在正文概括。

每位顾问的用途价值应说明，附本人签署合作信及简历；设备资源提供者等也需支持信。申请机构与合作机构、实验室的项目和行政安排，应有负责人致 PI 的承诺信。
\sourcepage{630}
\section{良好申请书的特征}
首先是研究策略的科学质量：好问题、周密可行的设计方法，以及有经验、技能、投入的团队。其次是表达清楚：简洁、有吸引力、组织良好、排版美观、没有错误，让人相信实际工作也会同样出色。

评审常面对厚厚一堆申请，即使快速阅读也应看出优点。\textbf{按具体目标组织，用有意义小标题划分短节，以图表打断长文，帮助理解。}原书所述 NIH 指南建议段首先用粗体主题句，忙碌者扫读即可把握要点。既提供足够细节说服专家，也吸引更大数量、不太熟悉领域的评审。

大多数评审厌恶夸大和过度包装。夸大重要性或能力会引起怀疑；可以热情，却应现实承认局限，因为评审善于识别设计与可行性问题。
\sourcepage{631}
近完成但仍有时间修改时，请未参与写作、熟悉领域的同事评论，也可请写作高手改进风格和清晰度。提交前务必读纸质稿，不要只依赖拼写语法软件。

\section{寻找研究支持}
也应寻找无需正式申请的机会，例如分析他人数据，或从资深研究者、部门取得小额支持。这样更快简单，但范围有限；学术晋升也常部分取决于争取经费的成绩。

医学研究经费主要有四类：
\begin{itemize}
\item 政府，主要 NIH，也包括退伍军人事务部、CDC、医疗保健研究与质量局、国防部及联邦、州、县机构；
\item 基金会、专业学会（如美国心脏协会、美国癌症协会）及个人捐赠；
\item 企业，尤其药品和器械公司；
\item 机构内部资源，如所在大学。
\end{itemize}
过程复杂、竞争激烈，有利于有经验且坚韧者；新研究者宜找这样的导师。下面介绍几类主要来源。

\subsection{NIH 基金与合同}
\sourcepage{632}
\term{R类资助}（R awards，包括 R01 和较小的 R03、R21）支持研究者自选问题或回应研究所公告的项目（\url{www.nimh.nih.gov/research-funding/grants/research-grants-r.shtml}）。\term{K类资助}（K awards，包括 K23、K01、K08 和本地 K12、KL2）为年轻研究者培训和职业发展提供工资与适量研究支持（\url{https://researchtraining.nih.gov/programs/career-development}）。多个评分部分影响总体评价（表20.3）。

\begingroup\small
\begin{longtable}{L{.22\textwidth}L{.12\textwidth}L{.57\textwidth}}
\caption*{\textbf{表20.3\quad 职业发展K类申请的评分部分}}\\\toprule
部分 & 建议篇幅 & 评审考虑\\\midrule\endfirsthead\toprule 部分 & 建议篇幅 & 评审考虑\\\midrule\endhead
申请人 & 1页 & 有独立研究潜力吗？既往培训经验适当吗？是否投入研究计划？有论文、摘要、经费等产出吗？\\
职业发展计划与目标 & 2--3页 & 能显著帮助走向独立吗？培训内容、范围、阶段及期限适合吗？\\
研究计划 & 7--8页 & 方法有科学和技术价值吗？符合当前阶段，有助获得培训所需技能吗？与职业目标相关吗？\\
\sourcepage{633}导师／导师组 & 0.5--1页 & 导师有资格经验提供所需训练吗？技能互补吗？培训角色说明充分吗？\\
机构承诺 & 1页 & 明确保证至少规定的投入用于研究培训吗？承诺充分吗？设施和机会适当、充足吗？\\\bottomrule
\end{longtable}\endgroup

研究所发起的申请，旨在推动 NIH 顾问委员会指定领域的研究，可采用\term{征求研究建议书}（Requests for Proposals，RFP）或\term{征求基金申请}（Requests for Applications，RFA）的形式。回应 RFP 时，研究者签订合同，完成 NIH 确定的具体研究活动；回应 RFA 时，主题领域由 NIH 规定，但研究者自行选择具体问题和研究设计。RFP 采用合同机制，向研究者所在机构补偿实现计划目标所需的成本；RFA 则采用基金机制，支持内容较开放的研究活动。

提交后依次经 NIH 行政初审、科学同行评审、研究所顾问委员会建议和所长最终决定。基金通常由全国各机构专家组成、各具专长的\term{研究评审组}（study sections）评审，名单和成员见 NIH 网站。
\sourcepage{634}
NIH 的评审与资助程序见\url{https://grants.nih.gov/grants/referral-and-review.htm}。申请提交后，由科学评审中心（Center for Scientific Review，CSR）分至某个评审组（图20.2），再指定一名主要评审者和两名或更多次要评审者。每人分别就意义、创新、研究方法、研究者和环境按1--9分评分，并对研究可能产生的总体影响作综合评价。总体影响得1分表示申请极强，基本没有弱点；9分表示存在严重弱点、优点很少。

指定评审者的评分先向评审组公布，分数较好的半数申请由全体委员讨论，其余“分流”而不讨论；少数申请因某些不清楚之处需要澄清，延至4个月后的下一轮。讨论后，指定评审者再次提出评分，分数可能因讨论而改变；随后全体委员无记名评分。平均分乘以10，得到10分（最好）至90分（最差）的总体评分。对于 R01，NIH 还把总体影响分换为百分位，表示该申请相对于同一评审组过去三次会议所评申请的排序。各研究所依据最终影响分或百分位，确定资助的优先次序。

\sourcepage{635}
\begin{figure}[p]\centering
\begin{tikzpicture}[font=\small,>=Stealth,x=1cm,y=1cm]
\node[font=\bfseries,align=center] at (3.3,0){NIH 的方式};
\node[font=\bfseries,align=center] at (11.2,0){基金会的典型方式};
\node (ngov) at (3.3,-.8){美国政府};
\node[align=center] (nfund) at (3.3,-2.8){NIH 经费分配至各研究所};
\draw[->] (ngov)--(nfund);
\node[align=left,text width=2cm] at (4.7,-1.7){国会拨款};
\coordinate (nfork) at (3.3,-4.2);
\draw (nfund.south)--(nfork);
\node[align=left,text width=2.7cm] at (5.1,-3.8){顾问委员会确定优先事项};
\node[align=center,text width=2.8cm] (ni) at (1.6,-5.1){研究者发起};
\node[align=center,text width=2.8cm] (ns) at (5,-5.1){研究所发起};
\draw[->] (nfork)--(ni);\draw[->] (nfork)--(ns);
\node[align=center,text width=2.8cm] at (1.6,-6.5){\textbf{R类}\\个人研究基金申请};
\node[align=center,text width=2.8cm] at (1.6,-8.2){\textbf{K类}\\职业发展};
\node[align=center,text width=2.8cm] at (5,-6.5){\textbf{RFP}\\征求研究建议书\\（合同）};
\node[align=center,text width=2.8cm] at (5,-8.2){\textbf{RFA}\\征求基金申请\\（基金）};
\node (fdonor) at (11.2,-.8){富有捐赠者};
\node (ffund) at (11.2,-2.8){基金会经费};
\draw[->] (fdonor)--(ffund);
\node[align=left,text width=2cm] at (12.6,-1.7){捐赠者赠款};
\coordinate (ffork) at (11.2,-4.2);\draw (ffund.south)--(ffork);
\node[align=left,text width=2cm] at (12.6,-3.8){董事会决定};
\node[align=center,text width=2.8cm] (fi) at (9.5,-5.1){研究者发起};
\node[align=center,text width=2.8cm] (fs) at (12.9,-5.1){基金会发起};
\draw[->] (ffork)--(fi);\draw[->] (ffork)--(fs);
\node[align=center,text width=2.8cm] at (9.5,-6.7){基金申请};
\node[align=center,text width=2.8cm] at (12.9,-6.7){征求基金申请};
\end{tikzpicture}
\caption{NIH 与基金会的经费来源和机制概览。}
\end{figure}

应在资深同事建议下预先选择合适评审组。各组不仅主题不同，专长和竞争申请质量也差异很大。可提出最多三个希望或不希望分到的组，但不保证遵从。

CSR 同时把申请分至 NIH 某研究所或中心，也可预先表达偏好。研究所按评分优先次序、顾问委员会审查和有时所长调整决定资助（图20.3）。
\sourcepage{636}
尚未获得 NIH 研究资助的职业早期新研究者，与资深研究者相比，通常适用较有利的资助百分位界限。某份申请若受到多个研究所关注，这些研究所有时会安排共同资助。

\begin{figure}[p]\centering
\begin{tikzpicture}[font=\small,>=Stealth,x=1cm,y=1cm]
\node[font=\bfseries] at (2.8,0){NIH 的方式};
\node[font=\bfseries] at (11.3,0){基金会的典型方式};
\node[align=center] (na) at (2.8,-1){R01基金申请};
\node[align=center,text width=5cm] (nc) at (2.8,-3.3){NIH 科学评审中心把 R01 分给：\\（1）一个评审组；\\（2）一个 NIH 研究所};
\node[align=center,text width=5cm] (nr) at (2.8,-6.3){研究所据此优先次序，\\将各评审组收到的全部申请排序};
\node[align=center,text width=5cm] (nf) at (2.8,-9.2){按优先次序资助申请};
\draw[->] (na)--(nc);\draw[->] (nc)--(nr);\draw[->] (nr)--(nf);
\node[align=left,text width=2.8cm] at (5.7,-4.9){评审组确定优先次序};
\node[align=left,text width=2.8cm] at (5.7,-7.8){顾问委员会审查次序，所长作最终决定};
\node[align=center] (fa) at (11.3,-1){基金申请};
\node[align=center,text width=4cm] (ff) at (11.3,-9.2){董事会决定哪些项目获资助};
\draw[->] (fa)--(ff);
\node[align=left,text width=3cm] at (13.1,-5){基金会工作人员提出是否资助的建议};
\end{tikzpicture}
\caption{NIH 与基金会评审申请的程序。}
\end{figure}

评审后，研究者收到\term{评审摘要}（summary statement），包括评分及委员会详细评论批评。首次申请常未获资助，NIH 允许修订重交一次；若评分评论提示可以改善，重交可能成功，但若缺乏创新意义，较难提高热情。
\sourcepage{637}
相关项目官员参加会议，宜与其讨论，因为书面意见常会前拟定，未必涵盖全部讨论。

无需照做每项建议，但应尽可能修订满足批评，并解释不修改的理由。NIH 限一页“引言”说明变化；可粗体或斜体概括主要批评，再简述相应修改。也可作为与旧评分无关的新申请反复提交，但应更新，可能还需修订；不确定应问 NIH 或查规定（\url{https://grants.nih.gov/grants/policy/resubmission_q&a.htm}）。

\subsection{基金会和专业学会}
多数私人基金会，如 Robert Wood Johnson Foundation，只支持特定方向。疾病基金会和心脏、癌症等学会也设研究项目，许多面向年轻人。总额远少于 NIH，往往希望支持 NIH 不太可能资助的主题或方法；少数有医疗质量等领域职业发展奖励。Foundation Center（\url{http://candid.org/}）提供可搜索名录、联系信息和写作建议。
\sourcepage{638}
各基金会程序不同，通常较快，常由工作人员建议、董事会核准的行政决策，而非同行评审。

应询问导师、查目标与近期资助网站，若适合则联系人员，介绍项目、判断兴趣并征求提交指导。许多要求先写信说明背景、主要目标、研究者资格、大致期限与成本；有兴趣再要详细申请。

\subsection{企业支持}
生产药物和器械的公司是重要的资助来源，尤其支持新治疗的随机试验。企业对临床研究最大的一类支持，是与临床研究场所签订合同，为检验新药或新器械的多中心试验招募参与者。这些大型试验通常由企业资助方运行，而且常通过与临床研究组织（clinical research organization，CRO）签约来实施；有时也由学术协调中心设计和管理。

通常按每名参与者固定付费，达到全研究目标即停止招募。人数足够可能收入超过成本，余额归机构，有时可用于其他研究；人数过少则亏损。参与前应确保机构行政与伦理能及时批准，来得及入足够人数。
\sourcepage{639}
企业尤其营销部门经费，常流向有助销售的主题。企业试验常由自家统计师分析、医学写作者起草，可能为产品宣传偏转信息（第7章）。参与署名者应警惕，确保分析周密、结果客观。理想情况下，指导或发表委员会应主要或全部由非资助方雇员的研究者组成。

有些公司接受研究者发起的小项目，如治疗机制效果或相关疾病流行病学，也可提供药物器械和匹配安慰剂。通常由研究者分析、撰稿。

\subsection{机构内部支持}
大学研究机构给本单位人员的金额通常较小，但取得更快（数周至数月）、成功率更高。可能限先导研究、争取外部经费或设备等用途，常专为年轻教师，是初学者领导获资助项目的机会。规定差异大，应先咨询行政人员。

\section{小结}
\sourcepage{640}
\begin{enumerate}
\item 申请书是用于请求经费的方案扩展，另含预算、行政与支持信息。
\item 新研究者先咨询资深同事，决定问题与机构，再读指南、联系科学管理人员。
\item 写作常比预期久：组队指定 PI、按指南列提纲、排书面任务时限、找范例、定期开会；请知情同事审阅，反复修订，最后细致润色。
\item 主要部分为摘要、预算及理由、简历与资源、至关重要的具体目标，以及含意义、创新和方法的策略，并介绍既往工作。
\item 好申请需好问题、计划和团队，也需清楚表达；逻辑提纲说明权衡，用小标题、图表突出优点，免被忙碌评审漏看。
\item 四类主要来源：
\begin{enumerate}
\item NIH 等政府是最大来源，同行和行政评审较慢，但支持广泛研究、合同与职业发展；
\item 基金会和学会常关注 NIH 常规范围外的有前景问题，较快但规则各异；
\sourcepage{641}
\item 药品器械企业主要支持新产品试验，部分有相关主题的研究者发起基金；
\item 机构内部小额经费较快、成功率较高，适合先导研究和初学者。
\end{enumerate}
\end{enumerate}

\clearpage
\chapterappendix{20A}{第20章习题：撰写研究经费申请书}
\sourcepage{642}
\begin{enumerate}
\item 具体目标之后，研究策略的三个主要部分是什么？
\item NIH 评审中，除环境外，还有哪四项按1--9评分并用于总体影响分？
\item 一项研究比较自动充气袖带、腕表与手机应用三种血压测量。全美超过50,000名20岁及以上志愿者参与在线心血管研究。研究者拟测各设备血压及变异，再通过病历随访新发和致死性卒中。
\begin{enumerate}
\item 提出两三个具体目标，每个一两句。
\item 写简短创新部分，说明至少两个新颖特点。
\end{enumerate}
\item 除 NIH 外，至少列出三个研究经费来源，最好适用于你的领域。
\end{enumerate}
